\documentclass[12pt]{article}
\usepackage[utf8]{inputenc}
\usepackage[T1]{fontenc}
\usepackage{hyperref}
\title{Ngazidja Comorian Language Resources for K-12 Curriculum Development}
\author{Compiled by Claude Sonnet 5.5 (claude-sonnet-5-5) for Andr\'as Kornai}
\date{}
\begin{document}
\maketitle

\section{Introduction}
Ngazidja Comorian (Shingazidja; ISO 639-3: zdj; Glottocode ngaz1238) is a Bantu language of Grande Comore, in the Comoros. It is not considered endangered. The best-documented resources found are academic works by Michel Lafon, which are in French and listed only as catalogue entries. No school textbooks or primers turned up in the search results.

\section{Primary Educational Resources}
No primers, textbooks or school programmes were found in search results.

\section{Linguistic and Reference Materials}
\begin{itemize}
\item Michel Lafon, \emph{Le shingazidja: une langue bantu sous influence arabe} (1987 thesis, INALCO, Paris), the main reference grammar according to Glottolog; the catalogue record I found does not link the text itself. The Glottolog record for Lafon's 1991 French--Shingazidja lexicon is here: \url{https://glottolog.org/resource/reference/id/21877}
\item Other Glottolog bibliographic records that appeared in search results (not individually verified): \url{https://glottolog.org/resource/reference/id/weball:11561}, \url{https://glottolog.org/resource/reference/id/weball:22794} and \url{https://glottolog.org/resource/reference/id/weball:11555}
\item CRDO/Cocoon language record for zdj: \url{https://cocoon.huma-num.fr/exist/crdo/lexvo.xql?code=zdj}
\item Wikipedia overview of the Comorian languages: \url{https://en.wikipedia.org/wiki/Comorian_languages}
\item Verbix page on Ngazidja Comorian: \url{https://docs.verbix.com/Languages/ComorianNgazidja}
\item Lexilogos Comorian page, with a dictionary link and sample texts: \url{https://www.lexilogos.com/comorien_dictionnaire.htm}
\end{itemize}

\section{Religious and Community Resources}
No religious or community resources specific to zdj were found in this search.

\section{Digital and Audio Resources}
No audio or digital resources specific to zdj were found in this search, apart from the online pages above.

\section{Curriculum Development Context}
The scholarly foundation is solid but thin: Lafon's grammar and lexicons, plus phonological work on Shingazidja tone and accent by Cassimjee and Kisseberth. Nothing found is designed for classrooms. A curriculum would need new pedagogical material built on these descriptions, developed with Comorian educators. French and Arabic influence on the language and on schooling in the Comoros should be taken into account.
\end{document}